%======================================================================%
\section{The Void and the Mirror Operator}
\label{sec:void-mirror}
%======================================================================%

This section states the \emph{only} primitive axioms of SJET and gives the
functorial definition of $\Mirror$ on which all rung theorems depend.

%----------------------------------------------------------------------%
\subsection{The void}
\label{sec:void-mirror:void}
%----------------------------------------------------------------------%

\begin{definition}[Void]\label{def:void-mirror:void}
The \emph{void} $\Void$ is the trivial observable algebra: the two-element
Boolean algebra $\mathcal{B}_0=\{\bot,\top\}$ with $\bot<\top$, and its Stone
spectrum $S_0=\mathrm{Spec}(\mathcal{B}_0)$, a single point. The associated
division datum is the zero algebra $\{0\}$ of dimension $0$. No vertex, no
edge, no localisable quantum number exists at $\Void$.
\end{definition}

\begin{axiom}[Void]\label{ax:void-mirror:void}
Physical reality is anchored at the void $\Void$: the state with no
distinguishable observable. All structure arises by departure from $\Void$, not
by importing a graph or Lie group as primitive data.
\end{axiom}

%----------------------------------------------------------------------%
\subsection{Category of Stone rungs}
\label{sec:void-mirror:category}
%----------------------------------------------------------------------%

\begin{definition}[Stone rung category]\label{def:void-mirror:StoneCat}
Let $\mathbf{Stone}$ be the category of Stone spaces (compact, Hausdorff,
totally disconnected) with continuous maps. A \emph{rung datum} is a pair
$(S,D)$ with $S\in\mathbf{Stone}$ and $D$ a formally real observable algebra
at that rung. Morphisms preserve both $S$ and the algebraic rung label.
\end{definition}

\begin{definition}[Three mirror modes]\label{def:void-mirror:modes}
The mirror operator $\Mirror$ acts in three modes, triggered by rung index:
\begin{enumerate}[label=\textbf{M\arabic*.},leftmargin=*]
  \item \textbf{Discrete} ($n=0\to1$): append a Boolean generator $a_{n+1}$ to
    $\mathcal{B}_n$, doubling clopen sectors:
    $\mathcal{B}_{n+1}=\mathcal{B}_n[a_{n+1}]/(a_{n+1}^2=a_{n+1})$.
  \item \textbf{Arithmetic} ($1\to2\to3\to4$): Cayley--Dickson doubling
    $D_{n+1}=D_n\oplus D_n$ with product
    $(x,y)\cdot(u,v)=(xu-\bar v y,\,vx+y\bar u)$ and involution
    $\overline{(x,y)}=(\bar x,-y)$ (Appendix~\ref{sec:appendix:cayley}).
  \item \textbf{Jordan exit} ($4\to5$): because $\Mirror_{\mathrm{arith}}(\Oct)$
    is non-division (sedenions), the climb exits to
    $\jalg=H_3(\Oct)$, the unique formally real Jordan envelope of $\Oct$.
  \item \textbf{Exceptional} ($n\geq6$): Lie mirror
    $G\mapsto G\times G$ or Cartan involution on $G/H$
    (Section~\ref{sec:ladder:exceptional}).
\end{enumerate}
\end{definition}

\begin{definition}[Mirror operator]\label{def:void-mirror:mirror}
A \emph{mirror operator} is a endofunctor
$\Mirror:\mathbf{Stone}\to\mathbf{Stone}$ together with algebra maps such that:
\begin{enumerate}[label=(\roman*),leftmargin=*]
  \item $\Mirror$ is \textbf{involutive} on stabilized rungs:
    $\Mirror^2\cong\mathrm{id}$ (Boolean complement, field conjugation, coset
    exchange $\mathbf{16}\leftrightarrow\overline{\mathbf{16}}$).
  \item $\Mirror$ is \textbf{nontrivial} at $\Void$:
    $S_1=\Mirror(S_0)$ has $|S_1|=2$ clopen points.
  \item $\Mirror$ is \textbf{functorial}: compatible maps $f:S\to T$ induce
    $\Mirror(f):\Mirror(S)\to\Mirror(T)$.
\end{enumerate}
\end{definition}

\begin{axiom}[Mirror]\label{ax:void-mirror:mirror}
There exists $\Mirror$ as in Definition~\ref{def:void-mirror:mirror} whose
orbit from $\Void$ generates the unified ladder \eqref{eq:ladder:unified}. At
each rung, capacity balancing $\mathrm{Bal}_\capacity$ is the \emph{dual
mirror} (Theorem~\ref{thm:ladder:concave}). The climb is closed:
$\Mirror^{\mathrm{tot}}(\cdots(\E{6})\cdots)=\Void$.
\end{axiom}

%----------------------------------------------------------------------%
\subsection{Rung table and the Hurwitz--Jordan theorems}
\label{sec:void-mirror:rungs}
%----------------------------------------------------------------------%

\begin{definition}[Mirror power]\label{def:void-mirror:power}
Define rung algebras $A_n$ and dimensions $d_n$ by
\begin{equation}\label{eq:void-mirror:power}
  (S_n,A_n)\;=\;\Mirror^n(S_0,A_0),
  \qquad (S_0,A_0)=(S(\Void),\{0\}),
  \qquad d_n=\dim_{\R} A_n.
\end{equation}
\end{definition}

\begin{theorem}[Mirror rung table]\label{thm:void-mirror:table}
The first six arithmetic rungs of $\Mirror^n(\Void)$ are:
\begin{center}
\renewcommand{\arraystretch}{1.2}
\begin{tabular}{@{}rclcl@{}}
\toprule
$n$ & $A_n$ & $d_n$ & Stone points $|S_n|$ & Mirror mode \\
\midrule
$0$ & $\{0\}$ & $0$ & $1$ & void \\
$1$ & $\R$ & $1$ & $2$ & discrete \\
$2$ & $\C$ & $2$ & $4$ & arithmetic \\
$3$ & $\Quat$ & $4$ & $8$ & arithmetic \\
$4$ & $\Oct$ & $8$ & $16$ & arithmetic (Hurwitz terminal) \\
$5$ & $\jalg$ & $27$ & $27$ & Jordan exit \\
\bottomrule
\end{tabular}
\end{center}
In particular $\Mirror^4(\Void)$ carries the octonions and
$\Mirror^5(\Void)\cong\jalg$.
\end{theorem}

\begin{proof}
\emph{Discrete step} ($n=0\to1$): adjoining one Boolean atom to $\mathcal{B}_0$
gives the four-element algebra $\{\bot,a,\neg a,\top\}$ with Stone spectrum
two points; this is the first observable distinction, identified with ordered
data $\R$ ($d_1=1$).

\emph{Arithmetic steps} ($n=1,2,3$): Cayley--Dickson doubling doubles
dimension: $1\to2\to4\to8$. Hurwitz's theorem
(Theorem~\ref{thm:ladder:hurwitz}) proves no further \emph{normed division}
algebra exists at $d=16$.

\emph{Jordan exit} ($n=4\to5$): $\Mirror_{\mathrm{arith}}(\Oct)$ is the
sedenion algebra $\mathbb{S}$ ($d=16$), which fails division
(Lemma~\ref{lem:void-mirror:sedenion}). The climb therefore exits to the
unique formally real Jordan algebra $H_3(\Oct)$ with $d_5=27$ and cubic norm
$N(X)=\det X$ (Proposition~\ref{prop:ladder:terminal}). The Stone spectrum
$S_5$ has $27$ clopen cells matching the $27$ matrix coordinates.
\end{proof}

\begin{lemma}[Sedenion failure]\label{lem:void-mirror:sedenion}
Cayley--Dickson doubling of $\Oct$ produces $\mathbb{S}$ of dimension $16$ with
zero divisors; $\mathbb{S}$ is not a division algebra. Hence
$\Mirror^5_{\mathrm{arith}}(\Void)\neq\mathbb{S}$ in the division category.
\end{lemma}

\begin{corollary}[Hurwitz terminal at $n=4$]\label{cor:void-mirror:hurwitz}
$\Mirror^4(\Void)$ is the \emph{last} normed division rung. The next stable
mirror power is $\Mirror^5(\Void)\cong\jalg$, not a $16$-dimensional division
algebra.
\end{corollary}

%----------------------------------------------------------------------%
\subsection{The climb operator}
\label{sec:void-mirror:climb-op}
%----------------------------------------------------------------------%

\begin{definition}[Climb operator]\label{def:void-mirror:Climb}
The \emph{climb operator} is
\begin{equation}\label{eq:void-mirror:Climb}
  \Climb\;=\;\mathrm{Bal}_\capacity\,\circ\,\Mirror.
\end{equation}
The \emph{universe datum} is the profinite fixed point
\begin{equation}\label{eq:void-mirror:universe}
  \mathfrak{U}^\star\;=\;\varprojlim_{n\to\infty}\Climb^n(\Void).
\end{equation}
At $n=5$, $\Climb^5(\Void)$ stabilizes on the Albert--Cayley graph
(Definition~\ref{def:foundations:albert}) with $1|10|16$ capacity partition.
\end{definition}

\begin{proposition}[Involution law]\label{prop:void-mirror:involution}
On each stabilized rung, $\Mirror^2\cong\mathrm{id}$ up to gauge:
Boolean complement, field conjugation, and coset exchange
$\mathbf{16}_{-3}\leftrightarrow\overline{\mathbf{16}}_{+3}$ are involutions.
\end{proposition}